%====================================================================
% 厚生労働省推奨 履歴書 — モハメド フアド
%====================================================================
\documentclass[9.5pt, a4paper]{article}
\usepackage{fontspec}
\usepackage{xeCJK}
\usepackage{geometry}
\usepackage{xcolor}
\usepackage[bookmarks=false]{hyperref}
\usepackage{tabularx}
\usepackage{array}
\usepackage{multirow}
\usepackage{colortbl}
\setCJKmainfont{ipaexm.ttf}[BoldFont={ipaexm.ttf}]
\setCJKsansfont{ipaexm.ttf}
\setmainfont{Helvetica Neue}
\geometry{a4paper,top=1.0cm,bottom=1.0cm,left=1.5cm,right=1.5cm}
\hypersetup{hidelinks}
\pagestyle{empty}
\setlength{\parindent}{0pt}
\renewcommand{\arraystretch}{1.3}

\definecolor{gridcolor2}{RGB}{170, 185, 200}
\definecolor{bannercolor}{RGB}{240, 244, 248}

\arrayrulecolor{gridcolor2}

\begin{document}
\begin{center}
  {\Large\bfseries\textcolor{darkgray}{履 \quad 歴 \quad 書}}\\[4pt]
  \hfill {\small\textcolor{gray}{2026年6月27日 現在}}
\end{center}

\vspace{4pt}
\begin{tabularx}{\linewidth}{|p{2cm}|X|p{1.8cm}|p{3.0cm}|p{3.0cm}|}
\hline
\rowcolor{bannercolor}
ふりがな & \multicolumn{3}{l|}{もはめど ふあど} & \multirow{4}{*}{\parbox[c][1.8cm][c]{2.6cm}{\centering\footnotesize\textcolor{gray}{写真貼付}\\[3pt]\tiny\textcolor{gray}{(縦4cm $\times$ 横3cm)}}} \\ \cline{1-4}
氏名 & \multicolumn{3}{l|}{\large\textbf{モハメド フアド}} & \\ \cline{1-4}
生年月日 & \multicolumn{3}{l|}{2004年2月28日 \quad （満 22 歳）} & \\ \cline{1-4}
性別（任意） &  & 電話番号 & 080-7535-2988 & \\ \hline
\rowcolor{bannercolor}
現住所 & \multicolumn{4}{l|}{東京都世田谷区} \\ \hline
メール & \multicolumn{4}{l|}{\href{mailto:mohamed.fuad.jp@gmail.com}{mohamed.fuad.jp@gmail.com}} \\ \hline
\end{tabularx}

\vspace{8pt}
\begin{tabularx}{\linewidth}{|p{1.8cm}|p{1.0cm}|X|}
\hline
\rowcolor{bannercolor}
\multicolumn{1}{|c|}{\textbf{年}} & \multicolumn{1}{c|}{\textbf{月}} & \multicolumn{1}{c|}{\textbf{学歴・職歴}} \\ \hline
 &  & \hfill \textbf{学\quad 歴} \hfill \\ \hline
2024 & 4 & 東海大学 情報通信学部 情報通信学科 入学 \\ \hline
2028 & 3 & 東海大学 情報通信学部 情報通信学科 卒業予定 \\ \hline
 &  & \hfill \textbf{職\quad 歴} \hfill \\ \hline
2023 & 2 & 日本航空株式会社 入社 \\ \hline
2023 & 4 & 日本航空株式会社 退社 \\ \hline
2023 & 4 & ホテルSUI赤坂 入社 \\ \hline
2023 & 6 & アルティウスリンク株式会社 入社 \\ \hline
2023 & 7 & ホテルSUI赤坂 退社 \\ \hline
 &  & 現在に至る \\ \hline
 &  & \hfill 以上 \\ \hline
 &  &  \\ \hline
\end{tabularx}

\vspace{8pt}
\begin{tabularx}{\linewidth}{|p{1.8cm}|p{1.0cm}|X|}
\hline
\rowcolor{bannercolor}
\multicolumn{1}{|c|}{\textbf{年}} & \multicolumn{1}{c|}{\textbf{月}} & \multicolumn{1}{c|}{\textbf{資格・免許}} \\ \hline
2024 & 7 & 日本語能力試験 N2 合格 \\ \hline
 &  &  \\ \hline
\end{tabularx}

\vspace{8pt}
\begin{tabularx}{\linewidth}{|X|}
\hline
\rowcolor{bannercolor} \textbf{自己PR・志望動機・アピールポイント} \\ \hline
\parbox[t][5.3cm][t]{\dimexpr\linewidth-2\tabcolsep}{
  \vspace{6pt}
  \fontsize{9.5}{14}\selectfont
  私の強みは、利用者の課題を整理し、必要な技術を学びながら動く形まで実装し切る力です。Tutor-Systemでは、PDF閲覧、出典を確認できるAI対話、音声学習、復習機能を一つの学習体験として設計しました。WebDropでは、ブラウザAPIの制約を検証しながら、WebRTCによる端末間転送と近接確認機能を実装しました。また、翻訳業務では、正確さだけでなく、相手や状況に応じて伝わる表現を選ぶことを大切にしています。開発力と日英両言語でのコミュニケーション力を活かし、チームの中で着実に価値を届けます。
  
  \vspace{8pt}
  \textbf{技術スキル:}
  \\ \textbf{言語:} TypeScript, JavaScript, Python, Swift, SQL, Java, C \quad \textbf{フレームワーク:} React, Node.js, Express, Tailwind CSS, SwiftUI, PWA
  \\ \textbf{開発環境・基盤:} Git, GitHub, AWS (Amplify/Cognito/SES), Docker, SQLite, MySQL, WebRTC
} \\ \hline
\end{tabularx}

\end{document}
